% ===========================================================================
%  CAPITULO 4 — Eixo 2: Experiencias (Bundles Estaticos, CDN, Assets e Planos)
% ===========================================================================
\chapter{Eixo 2 --- Experiências (Bundles Estáticos, CDN e Armazenamento)}
\label{cap:eixo-experiencias}

Na Plataforma Unooria, uma regra arquitetural inviolável determina que
\textbf{experiências pedagógicas (\arq{apps/experiencias/*}) nunca viram
serviços de contêiner no Cloud Run}. Elas são aplicações cliente puras (Vite +
React + \cd{@edu/sdk}), compiladas para arquivos estáticos e publicadas em
armazenamento de objetos atrás de CDN sob caminhos versionados e imutáveis:
\cd{/\{expId\}/\{versao\}/}. Toda inteligência, correção de nota, estado de
sessão e chamada de IA acontece no servidor (\arq{apps/api}).

Além do bucket dos \emph{bundles} das experiências (\cd{-exp}), este eixo
abrange os três outros buckets de objetos que sustentam o funcionamento das
missões e das ferramentas do professor (\arq{storage.tf}): \cd{-assets} (mídia
pesada pública e imagens geradas por IA), \cd{-planos} (planos de aula em
Markdown e áudios DUA dos professores, privados e versionados) e
\cd{-exp-mapas} (\emph{sourcemaps} privados para desofuscação de pilhas de erro).

\begin{fichaeixo}{Ficha do Eixo 2 --- Experiências e Armazenamento de Objetos}
\begin{campos}
  \item[Arquivos de referência] \arq{storage.tf}, \arq{loadbalancer.tf},
    \arq{scripts/publicar-manifesto.mjs} e \arq{apps/experiencias/AGENTS.md}.
  \item[Aderência por nuvem] \sgcp~\nativo \quad \saws~\nativo \quad
    \sazure~\nativo
  \item[Ponto focal] Este é o eixo de \textbf{maior paridade e menor atrito de
    migração} entre as três nuvens. GCS + Cloud CDN, Amazon S3 + CloudFront e
    Azure Blob Storage + Azure Front Door oferecem garantias de durabilidade
    (11 noves), versionamento de objetos, políticas de ciclo de vida e controle
    de cache por cabeçalho \cd{Cache-Control} virtualmente idênticas.
\end{campos}
\end{fichaeixo}

\section{Bundles estáticos imutáveis, CDN e isolamento de origem (\texttt{-exp})}

O contrato entre o Hub e a experiência embutida no \cd{<iframe>} depende de duas
propriedades da infraestrutura de entrega estática:

\begin{enumerate}
  \item \textbf{Imutabilidade e \emph{rollback} instantâneo:} publicar uma
    experiência copia os arquivos com \emph{hash} para
    \cd{/\{expId\}/\{versao\}/assets/*} com \cd{Cache-Control: public,
    max-age=31536000, immutable}, e o \arq{index.html} + \arq{experiencia.json}
    com \cd{Cache-Control: public, max-age=60}. Ativar ou reverter uma versão
    não apaga nem move arquivos no bucket --- apenas aponta o manifesto para a
    versão desejada (\arq{publicar-manifesto.mjs}), fazendo o \emph{rollback}
    propagar em até 60 segundos sem invalidar o cache global da CDN.
  \item \textbf{Origem dedicada para o \cd{postMessage}:} o Hub envia o token
    JWT para o \cd{<iframe>} via \cd{postMessage} com \cd{targetOrigin}
    explícito (nunca \cd{*}). Se as experiências fossem servidas por
    \cd{https://storage.googleapis.com/\{bucket\}} (ou \cd{s3.amazonaws.com}),
    elas dividiriam a mesma origem web (\emph{origin}) com qualquer outro bucket
    público daquela nuvem. Por isso o Load Balancer/CDN expõe o bucket sob
    domínio próprio do contrato (\cd{exp\_domain}, ex.:
    \cd{https://exp-entreredes.unooria.com}).
\end{enumerate}

\begin{pecaarq}{EXP-01}{Entrega de Bundles Estáticos e Manifesto por CDN (\texttt{\{contrato\}-exp})}{\sgcp~\nativo\enspace\saws~\nativo\enspace\sazure~\nativo}
\begin{campos}
  \item[Como é hoje (GCP)] Bucket \textbf{Cloud Storage}
    (\cd{\{project\_id\}-exp}) com \emph{Uniform Bucket-Level Access}, conectado
    ao Global External Load Balancer como \cd{google\_compute\_backend\_bucket}
    com \textbf{Cloud CDN} ativado no modo \cd{USE\_ORIGIN\_HEADERS}
    (\arq{loadbalancer.tf}). A publicação no CI usa \cd{gcloud storage rsync}
    com cabeçalhos diferenciados para \arq{assets/} versus \arq{index.html} e
    \arq{experiencia.json}.
  \item[Equivalente AWS] Bucket \textbf{Amazon S3} privado (com \emph{S3 Block
    Public Access} ligado em 100\% do bucket) servido exclusivamente pelo
    \textbf{Amazon CloudFront} usando \textbf{Origin Access Control (OAC)}. É um
    pequeno ganho de segurança sobre o Backend Bucket clássico do GCP: na AWS, o
    bucket S3 nem precisa ter leitura pública (\cd{allUsers}), pois apenas a
    distribuição CloudFront assinada via SigV4 consegue ler os objetos. A
    Política de Cache do CloudFront (\cd{UseOriginCacheControlHeaders}) respeita
    exatamente o \cd{max-age=60} do \arq{index.html} e o \cd{immutable} dos
    \emph{chunks} JS/CSS, e o deploy usa \cd{aws s3 sync}.
  \item[Equivalente Azure] Conta de armazenamento \textbf{Azure Blob Storage}
    (contêiner \cd{exp}) atrás do \textbf{Azure Front Door} com CDN habilitado e
    respeito aos cabeçalhos \cd{Cache-Control} da origem. No Front Door
    \emph{Premium}, o contêiner Blob também pode permanecer 100\% privado via
    \emph{Private Link} entre a borda da Microsoft e o Storage Account; no
    \emph{Standard}, usa acesso de leitura de blobs anônimos restrito ao
    contêiner, igual ao GCS. O deploy usa \cd{az storage blob sync} ou \cd{azcopy
    sync}.
\end{campos}
\end{pecaarq}

\section{Mídia pesada e imagens geradas por IA (\texttt{-assets})}

Para manter o repositório Git enxuto (regra inviolável: nenhum binário $>2\,\text{MB}$
no Git), vídeos de abertura das missões (como os clipes MP4 da \cd{exp02em1b}),
modelos 3D (\arq{mascote-unooria.glb}), áudios longos de podcasts e todas as
imagens \arq{.webp} geradas pela IA no preparo do professor (deduplicadas por
SHA-256 da descrição) moram no bucket \cd{\{contrato\}-assets}.

\begin{pecaarq}{EXP-02}{Bucket Público de Assets e Escrita pela API (\texttt{\{contrato\}-assets})}{\sgcp~\nativo\enspace\saws~\nativo\enspace\sazure~\nativo}
\begin{campos}
  \item[Como é hoje (GCP)] Bucket \textbf{Cloud Storage}
    (\cd{\{project\_id\}-assets}) exposto sob \cd{assets\_domain} no mesmo Load
    Balancer + Cloud CDN, com política CORS permitindo \cd{GET}/\cd{HEAD} a
    partir das origens do Hub e das Experiências (necessário para carregamento
    de texturas WebGL/Three.js no mascote 3D e \cd{<canvas>}). A conta de
    serviço da API possui \cd{roles/storage.objectAdmin} apenas neste bucket e
    no de planos, gravando as imagens convertidas em \arq{.webp} via \cd{sharp}
    com \cd{Cache-Control: public, max-age=31536000, immutable}.
  \item[Equivalente AWS] Bucket \textbf{Amazon S3} (\cd{\{contrato\}-assets})
    atrás do CloudFront (\cd{assets\_domain}) com \emph{Response Headers Policy}
    para CORS (\cd{Access-Control-Allow-Origin} dinâmico por lista de origens
    do contrato). A \emph{Task Role} da API no ECS Fargate recebe permissão IAM
    \cd{s3:PutObject} e \cd{s3:HeadObject} restrita ao ARN do bucket, usando o
    \cd{@aws-sdk/client-s3} (ou a interface compatível S3 do próprio GCS, se
    desejado unificar código).
  \item[Equivalente Azure] Contêiner \cd{assets} no \textbf{Azure Blob Storage}
    atrás do Azure Front Door, com regras nativas de CORS no serviço de Blob e
    permissão RBAC \emph{Storage Blob Data Contributor} atribuída à
    \emph{Managed Identity} do Container App da API.
\end{campos}
\end{pecaarq}

\section{Guarda versionada de planos de aula e sourcemaps privados}

Os dois últimos buckets do Eixo 2 são estritamente privados (\cd{public\_read =
false}) e nunca são expostos no Load Balancer nem na CDN:

\begin{pecaarq}{EXP-03}{Buckets Privados: Planos Versionados (\texttt{-planos}) e Mapas (\texttt{-exp-mapas})}{\sgcp~\nativo\enspace\saws~\nativo\enspace\sazure~\nativo}
\begin{campos}
  \item[Planos e anexos DUA (\texttt{-planos})] Guarda os planos de aula anuais
    gerados pela IA e editados pelos professores, além dos áudios e materiais
    de Desenho Universal para a Aprendizagem (DUA). Possui \textbf{Object
    Versioning} ativo e regra de ciclo de vida que retém versões não atuais por
    \cd{planos\_historico\_dias} (180 dias em produção), garantindo que uma
    edição acidental de plano nunca perca o texto anterior. Só a API lê e
    escreve; o navegador nunca acessa este bucket diretamente.
  \item[Sourcemaps (\texttt{-exp-mapas})] Guarda os arquivos \arq{*.js.map}
    gerados no build das experiências (\cd{sourcemap: 'hidden'}), permitindo
    desofuscar um \emph{stack trace} capturado em produção sem expor o código
    original aos estudantes no bucket público \cd{-exp}.
  \item[Tradução AWS e Azure] Tanto o \textbf{Amazon S3} (\emph{S3 Versioning} +
    regra \cd{NoncurrentVersionExpiration} de 180 dias) quanto o \textbf{Azure
    Blob Storage} (\emph{Blob Versioning} + política de gerenciamento de ciclo
    de vida \cd{daysAfterVersionCreationGreaterThan: 180}) implementam esse
    comportamento de maneira declarativa no Terraform, com paridade funcional de
    100\%.
\end{campos}
\end{pecaarq}

\begin{table}[H]
\centering
\caption{Comparativo técnico dos 4 buckets do Eixo de Experiências nas três nuvens}
\label{tab:comparativo-buckets-exp}
\footnotesize
\begin{tabularx}{\linewidth}{@{}L{28mm}L{38mm}L{38mm}X@{}}
\toprule
\headrow \thd{Recurso / Requisito} & \thd{GCP (Cloud Storage)} & \thd{AWS (Amazon S3)} & \thd{Azure (Blob Storage)}\\
\midrule
SDK na API (\arq{apps/api}) &
  \cd{@google-cloud/storage} (ADC nativo) &
  \cd{@aws-sdk/client-s3} (IAM Task Role) &
  \cd{@azure/storage-blob} (Managed Identity)\\
\zebra Sincronização no CI &
  \cd{gcloud storage rsync} &
  \cd{aws s3 sync} &
  \cd{azcopy sync} / \cd{az storage blob sync}\\
Proteção da origem pública &
  Bucket público lido pelo Cloud CDN &
  Bucket privado lido via \textbf{CloudFront OAC} &
  Contêiner público ou privado via \textbf{Private Link}\\
\zebra Versionamento (\cd{-planos}) &
  \cd{versioning = true} + \cd{num\_newer\_versions} / \cd{days\_since\_noncurrent} &
  \cd{aws\_s3\_bucket\_versioning} + \cd{noncurrent\_version\_expiration} &
  \cd{versioning\_enabled = true} + \cd{ManagementPolicy} por idade da versão\\
Custo base (armazenamento Standard) &
  $\approx$ US\$\,0,020 / GB-mês (\cd{us-east1}) &
  $\approx$ US\$\,0,023 / GB-mês (\cd{us-east-1}) &
  $\approx$ US\$\,0,018 a 0,021 / GB-mês (\emph{Hot tier})\\
\bottomrule
\end{tabularx}
\end{table}
